\begin{center}[6.]\par\medskip\textit{On the Mathematical Theory of Electrodynamic Effects.}\end{center}
1.\quad \textsc{Ampère} expresses the mutual action of two current elements $\mathrm ds,\mathrm ds'$, with the current intensities denoted by $i,i'$, by the formula
\[
ii'(\sin\theta\sin\theta'\cos\omega+k\cos\theta\cos\theta')r^n\mathrm ds\,.\,\mathrm ds'
\]
He assumes that the action takes place along the straight line joining the elements. Positive and negative signs indicate attraction and repulsion, respectively. Here $r$ denotes the distance between the elements, $\theta$ and $\theta'$ the angles the elements $\mathrm ds,\mathrm ds'$ make with $r$, this line being taken in the same direction for both, and finally $\omega$ the angle between the planes through $r$ and $\mathrm ds$ on the one hand and $r$ and $\mathrm ds'$ on the other. From his experiments \textsc{Ampère} concluded that one must set $n=-2$, $k=-\frac{1}{2}$. The mutual attraction is therefore measured by
\[
\frac{ii'(\sin\theta\sin\theta'\cos\omega-\frac{1}{2}\cos\theta\cos\theta')\mathrm ds\,.\,\mathrm ds'}{rr}
\]
A negative value of this expression indicates repulsion.

2.\quad We also use $a,a'$ for the full lengths of the currents, $\rho$ for the square root of $r$, $x,y,z$ for the coordinates of any point in $s$, $x',y',z'$ for the coordinates of any point in $s'$, and $\varepsilon$ for the
\SourceLeaf{613}angle between two current elements $\mathrm ds,\mathrm ds'$. We distinguish partial differentiations with respect to $s$ and $s'$ by the symbols $\mathrm d,\mathrm d'$ and partial integrations by the signs $\int,\int'$. We have
\[
\begin{gathered}
\mathrm dr=2\rho\,\mathrm d\rho=\cos\theta\,.\,\mathrm ds,\qquad \mathrm d'r=2\rho\,\mathrm d'\rho=-\cos\theta'\,.\,\mathrm ds'\\
\rho^4=(x'-x)^2+(y'-y)^2+(z'-z)^2\\
2\rho^3\mathrm d\rho=-(x'-x)\mathrm dx-(y'-y)\mathrm dy-(z'-z)\mathrm dz\\
\begin{aligned}
6\rho\rho\,\mathrm d\rho\,.\,\mathrm d'\rho+2\rho^3\mathrm d\mathrm d'\rho
&=-\mathrm dx\,.\,\mathrm d'x'-\mathrm dy\,.\,\mathrm d'y'-\mathrm dz\,.\,\mathrm d'z'\\
&=-\cos\varepsilon\,.\,\mathrm ds\,.\,\mathrm ds'\\
2\rho^3\mathrm d\mathrm d'\rho&=(-\cos\varepsilon+\tfrac{3}{2}\cos\theta\,.\,\cos\theta')\mathrm ds\,.\,\mathrm ds'\\
&=-(\sin\theta\sin\theta'\,.\,\cos\omega-\tfrac{1}{2}\cos\theta\,.\,\cos\theta')\mathrm ds\,.\,\mathrm ds'
\end{aligned}
\end{gathered}
\]
\textsc{Ampère}’s expression above thus becomes
\[
-\frac{2ii'\mathrm d\mathrm d'\rho}{\rho}
\]

3.\quad We use the symbol $\delta$ before a quantity to denote its infinitesimal variations in so far as they depend on infinitesimal virtual displacements of points of the conductors $s,s'$. Virtual displacements are any imagined displacements compatible with the conditions imposed by the nature of the conductors and their external circumstances.

4.\quad The full action of the currents on each other is to be regarded as known if one can specify other forces acting at finitely or infinitely many individual points of the currents and equivalent to their forces, in the sense that the opposites of the specified forces balance them. Let $W$ be the sum of the latter forces, each multiplied by its virtual motion. Then, by the principle of virtual motions,
\[
W+\int\int'\delta r\,.\,\frac{2ii'\mathrm d\mathrm d'\rho}{\rho}=W+4ii'\int\int'\mathrm d\mathrm d'\rho\,.\,\delta\rho
\]
must be $=0$, or at least cannot be positive, for each system of virtual motions of the conductors that is compatible with the conditions in force. The current forces alone, on the other hand, hold the conductors in equilibrium if $\int\int'\delta\rho\frac{\mathrm d\mathrm d'\rho}{\rho}$ is $=0$ or negative for all virtual motions. The integrations here are to be taken from $s=0$ to $s=a$ and from $s'=0$ to $s'=a'$.

\SourceLeaf{614}5.\quad To see the nature of this integral, we expand the variation of $\int\int'\mathrm d\rho\,.\,\mathrm d'\rho$. We have
\[
\begin{aligned}
\delta\int\int'\mathrm d\rho\,.\,\mathrm d'\rho&=\int\int'\delta(\mathrm d\rho\,.\,\mathrm d'\rho)\\
&=\int\int'\delta\mathrm d\rho\,.\,\mathrm d'\rho+\int\int'\delta\mathrm d'\rho\,.\,\mathrm d\rho
\end{aligned}
\]
Now
\[
\mathrm d\int'\mathrm d'\rho\delta\rho=\int'\mathrm d\delta\rho\,.\,\mathrm d'\rho+\int'\delta\rho\,.\,\mathrm d\mathrm d'\rho
\]
or, on integrating from $s=0$ to $s=a$ and writing the values of $\int'\mathrm d'\rho\,.\,\delta\rho$ at $s=0$ and $s=a$ as $F',G'$,
\[
-F'+G'=\int\int'\mathrm d\delta\rho\,.\,\mathrm d'\rho+\int\int'\delta\rho\,.\,\mathrm d\mathrm d'\rho
\]
and likewise, writing the values of $\int\mathrm d\rho\,.\,\delta\rho$ at $s'=0$ and $s'=a'$ as $F,G$,
\[
-F+G=\int\int'\mathrm d'\delta\rho\,.\,\mathrm d\rho+\int\int'\delta\rho\,.\,\mathrm d\mathrm d'\rho
\]
Consequently
\[
\begin{aligned}
-F-F'+G+G'&=\int\int'\mathrm d\delta\rho\,.\,\mathrm d'\rho+\int\int'\mathrm d'\delta\rho\,.\,\mathrm d\rho\\
&\qquad+2\int\int'\delta\rho\,.\,\mathrm d\mathrm d'\rho\\
&=\delta\int\int'\mathrm d\rho\,.\,\mathrm d'\rho+2\int\int'\delta\rho\,.\,\mathrm d\mathrm d'\rho
\end{aligned}
\]

The virtual moment of the mutual action of the current sections $s,s'$ is therefore
\[
\begin{aligned}
V=4ii'\int\int'\mathrm d\mathrm d'\rho\,.\,\delta\rho&=-2ii'\delta\int\int'\mathrm d\rho\,.\,\mathrm d'\rho\\
&\qquad+2ii'(G+G'-F-F')
\end{aligned}
\]

6.\quad A current exerts a current-generating or electromotive force only when it arises or moves. The electromotive force of a current element $i\,\mathrm ds$ acts at every point with a magnitude inversely proportional to the distance $r$ and directly proportional to the current element projected onto this line. It acts along the line $r$ itself, but always in the opposite sense. It follows readily that the virtual moment of the electromotive force of the element $i\,\mathrm ds$ can be expressed by
\[
+\frac{i\,\mathrm ds}{r}\,.\,\frac{\mathrm dr}{\mathrm ds}\,.\,\delta r
\]
or by
\SourceLeaf{615}\[
4i\,\mathrm d\rho\,.\,\delta\rho
\]

In the conductor element $\mathrm ds'$ the electromotive force of the current element $i\,\mathrm ds$ is therefore
\[
=4i\,\mathrm d\rho\,.\,\frac{\mathrm d\rho}{\mathrm ds'}\,.\,\mathrm ds'=4i\,\mathrm d\rho\,.\,\mathrm d'\rho
\]

The double integral $4i\int\int'\mathrm d\rho\,.\,\mathrm d'\rho$ is therefore the total electromotive force exerted in the conductor section from $s'=0$ to $s'=a'$ by the current section $is$, from $s=0$ to $s=a$. We denote this force by $A$. Furthermore, the virtual moment of the electromotive force that the same current section exerts at a given point is
\[
=4i\int\mathrm d\rho\,.\,\delta\rho\qquad\text{from }s=0\text{ to }s=a
\]
If this at the beginning of the conductor section $s'$ is $=B$ and at its end $=C$, then
\[
B=4iF,\qquad C=4iG
\]
If $B',C'$ likewise denote the virtual moments of the electromotive force exerted by a current section $i's'$ from $s'=0$ to $s'=a'$ at the beginning and end points of $s$, then $B'=4i'F'$, $C'=4i'G'$.

We therefore have
\[
V=-\tfrac{1}{2}i'\delta A-\tfrac{1}{2}i'B+\tfrac{1}{2}i'C-\tfrac{1}{2}iB'+\tfrac{1}{2}iC'
\]

